% --- Archon physics formalization source begin ---
% archon:physics
% archon:covers PhyXMiniProblems/problem_phyx_mini_0656.lean
% archon:source-report reports/phyx_mini/problem_phyx_mini_0656.source.json

\chapter{Physics problem phyx\_mini\_0656}
\label{ch:PhyXMiniProblems_problem_phyx_mini_0656}

\paragraph{Problem source.}
\#\# Physical scenario

Figure shows the wave function of an electron in a rigid box. The electron energy is \$6.0\$ \$eV\$.

\#\# Question

How long is the box?

\#\# Answer choices

- A: A: 0.5 nm

- B: B: 0.75 nm

- C: C: 1.0 nm

- D: D: 2.0 nm

\#\# Figure caption (auxiliary; use the image as primary evidence)

The image is a graph depicting a wave function represented by a sinusoidal red curve. The graph has two axes: the vertical axis labeled \textbackslash{}( \textbackslash{}psi(x) \textbackslash{}) and the horizontal axis labeled \textbackslash{}( x \textbackslash{}). The red curve exhibits a sinusoidal pattern: it starts at a value on the vertical axis, oscillates above and below the horizontal axis, and then returns to a flat line parallel to the horizontal axis on both ends. The wave has a smooth and continuous shape with two full oscillations visible in the image.

\#\# Dataset metadata

Quantum Phenomena; Physical Model Grounding Reasoning, Spatial Relation Reasoning

\paragraph{Recorded answer/context.}
C: C: 1.0 nm

\paragraph{Figure/image path.}
/root/proposal\_for\_physic/science-mango/phyx\_mini\_run/phyx\_data/test\_image/656.png

\paragraph{Formalization target.}
create a compiling Lean file with sorry bodies at `PhyXMiniProblems/problem\_phyx\_mini\_0656.lean`. The Lean declarations must preserve the physical quantities, dimensions or dimensional roles, figure labels, governing-law hypotheses, and final relation expressed by this problem.
Use Mathlib/Physlib names found through LeanExplore where available. If a physics API is missing, introduce faithful local abstractions rather than scalar placeholder aliases.

\begin{theorem}[Physics formalization target]
\label{thm:physics:phyx_mini_0656:target}
\lean{PhyXMiniProblems.ProblemPhyXMini0656.problem_phyx_mini_0656}
\uses{def:physics:phyx-mini-0656:phyxminiproblems-problemphyxmini0656-lengthinnanometers, def:physics:phyx-mini-0656:phyxminiproblems-problemphyxmini0656-rigidboxelectronsetup, def:physics:phyx-mini-0656:phyxminiproblems-problemphyxmini0656-matchesrigidboxelectronscenario, def:physics:phyx-mini-0656:phyxminiproblems-problemphyxmini0656-matchesproblemenergyreadout, def:physics:phyx-mini-0656:phyxminiproblems-problemphyxmini0656-useselectronmassreferencedata, def:physics:phyx-mini-0656:phyxminiproblems-problemphyxmini0656-matchessuppliedwavefunctionfigure, def:physics:phyx-mini-0656:phyxminiproblems-problemphyxmini0656-representsrigidboxstationarymode, def:physics:phyx-mini-0656:phyxminiproblems-problemphyxmini0656-satisfiesinfiniterigidboxenergylaw, def:physics:phyx-mini-0656:phyxminiproblems-problemphyxmini0656-hasphysicalrigidboxparameters, def:physics:phyx-mini-0656:phyxminiproblems-problemphyxmini0656-recordeddatasetanswer, def:physics:phyx-mini-0656:phyxminiproblems-problemphyxmini0656-roundedtonearesttenth, def:physics:phyx-mini-0656:phyxminiproblems-problemphyxmini0656-isuniquematchingdisplayedlength}
The assigned autoformalize agent should translate this physics problem into Lean declarations in the covered file, with theorem and lemma proof bodies written as `by sorry`.
\end{theorem}
\begin{proof}
This is an autoformalization task, not a proof task. Produce statements that can later be proved by the physics prover without weakening the physical model.
\end{proof}

% --- Archon declaration topology begin ---
\section{Declaration topology}
\begin{definition}[Length Quantity]
  \label{def:physics:phyx-mini-0656:phyxminiproblems-problemphyxmini0656-lengthquantity}
  \lean{PhyXMiniProblems.ProblemPhyXMini0656.LengthQuantity}
  A nonnegative, unit-independent physical length
\end{definition}
\begin{proof}
  This is the declared model definition of Length Quantity.
\end{proof}
\begin{definition}[Mass Quantity]
  \label{def:physics:phyx-mini-0656:phyxminiproblems-problemphyxmini0656-massquantity}
  \lean{PhyXMiniProblems.ProblemPhyXMini0656.MassQuantity}
  A nonnegative, unit-independent physical mass
\end{definition}
\begin{proof}
  This is the declared model definition of Mass Quantity.
\end{proof}
\begin{definition}[length Readout]
  \label{def:physics:phyx-mini-0656:phyxminiproblems-problemphyxmini0656-lengthreadout}
  \lean{PhyXMiniProblems.ProblemPhyXMini0656.lengthReadout}
  \uses{def:physics:phyx-mini-0656:phyxminiproblems-problemphyxmini0656-lengthquantity}
  Read a physical length in a selected length unit
\end{definition}
\begin{proof}
  This is the declared model definition of length Readout.
\end{proof}
\begin{definition}[length In Meters]
  \label{def:physics:phyx-mini-0656:phyxminiproblems-problemphyxmini0656-lengthinmeters}
  \lean{PhyXMiniProblems.ProblemPhyXMini0656.lengthInMeters}
  \uses{def:physics:phyx-mini-0656:phyxminiproblems-problemphyxmini0656-lengthquantity, def:physics:phyx-mini-0656:phyxminiproblems-problemphyxmini0656-lengthreadout}
  SI metre readout of a physical length
\end{definition}
\begin{proof}
  This is the declared model definition of length In Meters.
\end{proof}
\begin{definition}[length In Nanometers]
  \label{def:physics:phyx-mini-0656:phyxminiproblems-problemphyxmini0656-lengthinnanometers}
  \lean{PhyXMiniProblems.ProblemPhyXMini0656.lengthInNanometers}
  \uses{def:physics:phyx-mini-0656:phyxminiproblems-problemphyxmini0656-lengthquantity, def:physics:phyx-mini-0656:phyxminiproblems-problemphyxmini0656-lengthreadout}
  Nanometre readout used by the displayed answer choices
\end{definition}
\begin{proof}
  This is the declared model definition of length In Nanometers.
\end{proof}
\begin{definition}[mass In Kilograms]
  \label{def:physics:phyx-mini-0656:phyxminiproblems-problemphyxmini0656-massinkilograms}
  \lean{PhyXMiniProblems.ProblemPhyXMini0656.massInKilograms}
  \uses{def:physics:phyx-mini-0656:phyxminiproblems-problemphyxmini0656-massquantity}
  SI kilogram readout of a physical mass
\end{definition}
\begin{proof}
  This is the declared model definition of mass In Kilograms.
\end{proof}
\begin{definition}[energy In Joules]
  \label{def:physics:phyx-mini-0656:phyxminiproblems-problemphyxmini0656-energyinjoules}
  \lean{PhyXMiniProblems.ProblemPhyXMini0656.energyInJoules}
  SI joule readout of a physical energy
\end{definition}
\begin{proof}
  This is the declared model definition of energy In Joules.
\end{proof}
\begin{definition}[energy In Electron Volts]
  \label{def:physics:phyx-mini-0656:phyxminiproblems-problemphyxmini0656-energyinelectronvolts}
  \lean{PhyXMiniProblems.ProblemPhyXMini0656.energyInElectronVolts}
  \uses{def:physics:phyx-mini-0656:phyxminiproblems-problemphyxmini0656-energyinjoules}
  Electron-volt readout, using Physlib's calibrated electron volt
\end{definition}
\begin{proof}
  This is the declared model definition of energy In Electron Volts.
\end{proof}
\begin{definition}[Particle Species]
  \label{def:physics:phyx-mini-0656:phyxminiproblems-problemphyxmini0656-particlespecies}
  \lean{PhyXMiniProblems.ProblemPhyXMini0656.ParticleSpecies}
  Species of the particle represented by the plotted state
\end{definition}
\begin{proof}
  This is the declared model definition of Particle Species.
\end{proof}
\begin{definition}[Confinement Model]
  \label{def:physics:phyx-mini-0656:phyxminiproblems-problemphyxmini0656-confinementmodel}
  \lean{PhyXMiniProblems.ProblemPhyXMini0656.ConfinementModel}
  Confining potential used in the physical model
\end{definition}
\begin{proof}
  This is the declared model definition of Confinement Model.
\end{proof}
\begin{definition}[Figure Label]
  \label{def:physics:phyx-mini-0656:phyxminiproblems-problemphyxmini0656-figurelabel}
  \lean{PhyXMiniProblems.ProblemPhyXMini0656.FigureLabel}
  Axis labels printed in the primary figure
\end{definition}
\begin{proof}
  This is the declared model definition of Figure Label.
\end{proof}
\begin{definition}[Lobe Sign]
  \label{def:physics:phyx-mini-0656:phyxminiproblems-problemphyxmini0656-lobesign}
  \lean{PhyXMiniProblems.ProblemPhyXMini0656.LobeSign}
  Sign of one consecutive lobe of the plotted real wave amplitude
\end{definition}
\begin{proof}
  This is the declared model definition of Lobe Sign.
\end{proof}
\begin{definition}[Rigid Box Wave Function Figure]
  \label{def:physics:phyx-mini-0656:phyxminiproblems-problemphyxmini0656-rigidboxwavefunctionfigure}
  \lean{PhyXMiniProblems.ProblemPhyXMini0656.RigidBoxWaveFunctionFigure}
  \uses{def:physics:phyx-mini-0656:phyxminiproblems-problemphyxmini0656-figurelabel, def:physics:phyx-mini-0656:phyxminiproblems-problemphyxmini0656-lobesign}
  Qualitative and integer-valued evidence read from the primary raster. The figure has no numerical axis scale and does not draw the vertical box walls; its flat exterior pieces lie on the horizontal axis
\end{definition}
\begin{proof}
  This is the declared model definition of Rigid Box Wave Function Figure.
\end{proof}
\begin{definition}[Rigid Box Electron Setup]
  \label{def:physics:phyx-mini-0656:phyxminiproblems-problemphyxmini0656-rigidboxelectronsetup}
  \lean{PhyXMiniProblems.ProblemPhyXMini0656.RigidBoxElectronSetup}
  \uses{def:physics:phyx-mini-0656:phyxminiproblems-problemphyxmini0656-lengthquantity, def:physics:phyx-mini-0656:phyxminiproblems-problemphyxmini0656-massquantity, def:physics:phyx-mini-0656:phyxminiproblems-problemphyxmini0656-particlespecies, def:physics:phyx-mini-0656:phyxminiproblems-problemphyxmini0656-confinementmodel, def:physics:phyx-mini-0656:phyxminiproblems-problemphyxmini0656-rigidboxwavefunctionfigure}
  Independent physical data for the electron state and box. In particular, boxLength is an observable field and is not defined from the recorded answer. The representative is retained because an L² state is an almost-everywhere equivalence class while the plotted sine curve is pointwise
\end{definition}
\begin{proof}
  This is the declared model definition of Rigid Box Electron Setup.
\end{proof}
\begin{definition}[Matches Rigid Box Electron Scenario]
  \label{def:physics:phyx-mini-0656:phyxminiproblems-problemphyxmini0656-matchesrigidboxelectronscenario}
  \lean{PhyXMiniProblems.ProblemPhyXMini0656.MatchesRigidBoxElectronScenario}
  \uses{def:physics:phyx-mini-0656:phyxminiproblems-problemphyxmini0656-rigidboxelectronsetup}
  Categorical facts stated in the physical scenario
\end{definition}
\begin{proof}
  This is the declared model definition of Matches Rigid Box Electron Scenario.
\end{proof}
\begin{definition}[Matches Problem Energy Readout]
  \label{def:physics:phyx-mini-0656:phyxminiproblems-problemphyxmini0656-matchesproblemenergyreadout}
  \lean{PhyXMiniProblems.ProblemPhyXMini0656.MatchesProblemEnergyReadout}
  \uses{def:physics:phyx-mini-0656:phyxminiproblems-problemphyxmini0656-energyinelectronvolts, def:physics:phyx-mini-0656:phyxminiproblems-problemphyxmini0656-rigidboxelectronsetup}
  The energy value explicitly supplied in the question
\end{definition}
\begin{proof}
  This is the declared model definition of Matches Problem Energy Readout.
\end{proof}
\begin{definition}[Uses Electron Mass Reference Data]
  \label{def:physics:phyx-mini-0656:phyxminiproblems-problemphyxmini0656-useselectronmassreferencedata}
  \lean{PhyXMiniProblems.ProblemPhyXMini0656.UsesElectronMassReferenceData}
  \uses{def:physics:phyx-mini-0656:phyxminiproblems-problemphyxmini0656-massinkilograms, def:physics:phyx-mini-0656:phyxminiproblems-problemphyxmini0656-rigidboxelectronsetup}
  Standard electron-mass data required to evaluate the energy law. This is reference data, separate from the problem and figure readouts
\end{definition}
\begin{proof}
  This is the declared model definition of Uses Electron Mass Reference Data.
\end{proof}
\begin{definition}[Matches Supplied Wave Function Figure]
  \label{def:physics:phyx-mini-0656:phyxminiproblems-problemphyxmini0656-matchessuppliedwavefunctionfigure}
  \lean{PhyXMiniProblems.ProblemPhyXMini0656.MatchesSuppliedWaveFunctionFigure}
  \uses{def:physics:phyx-mini-0656:phyxminiproblems-problemphyxmini0656-rigidboxwavefunctionfigure}
  Raw facts visible in the supplied wavefunction graph
\end{definition}
\begin{proof}
  This is the declared model definition of Matches Supplied Wave Function Figure.
\end{proof}
\begin{definition}[Represents Rigid Box Stationary Mode]
  \label{def:physics:phyx-mini-0656:phyxminiproblems-problemphyxmini0656-representsrigidboxstationarymode}
  \lean{PhyXMiniProblems.ProblemPhyXMini0656.RepresentsRigidBoxStationaryMode}
  \uses{def:physics:phyx-mini-0656:phyxminiproblems-problemphyxmini0656-lengthinmeters, def:physics:phyx-mini-0656:phyxminiproblems-problemphyxmini0656-rigidboxelectronsetup}
  Interpretation of the plot as the nth stationary infinite-well mode. The mode number counts half-wave lobes, and a pointwise representative has the usual sine profile inside [0,L] and vanishes outside. No numerical value of L is part of this governing relation
\end{definition}
\begin{proof}
  This is the declared model definition of Represents Rigid Box Stationary Mode.
\end{proof}
\begin{definition}[Satisfies Infinite Rigid Box Energy Law]
  \label{def:physics:phyx-mini-0656:phyxminiproblems-problemphyxmini0656-satisfiesinfiniterigidboxenergylaw}
  \lean{PhyXMiniProblems.ProblemPhyXMini0656.SatisfiesInfiniteRigidBoxEnergyLaw}
  \uses{def:physics:phyx-mini-0656:phyxminiproblems-problemphyxmini0656-lengthinmeters, def:physics:phyx-mini-0656:phyxminiproblems-problemphyxmini0656-massinkilograms, def:physics:phyx-mini-0656:phyxminiproblems-problemphyxmini0656-energyinjoules, def:physics:phyx-mini-0656:phyxminiproblems-problemphyxmini0656-rigidboxelectronsetup}
  The one-dimensional infinite-square-well spectrum Eₙ = (n π ℏ)² / (2 m L²). The equation uses coherent SI readouts. Physlib's Constants.ℏ is the reduced Planck constant in joule-seconds. This is a general law and does not assume the requested box length or any answer choice
\end{definition}
\begin{proof}
  This is the declared model definition of Satisfies Infinite Rigid Box Energy Law.
\end{proof}
\begin{definition}[Has Physical Rigid Box Parameters]
  \label{def:physics:phyx-mini-0656:phyxminiproblems-problemphyxmini0656-hasphysicalrigidboxparameters}
  \lean{PhyXMiniProblems.ProblemPhyXMini0656.HasPhysicalRigidBoxParameters}
  \uses{def:physics:phyx-mini-0656:phyxminiproblems-problemphyxmini0656-lengthinmeters, def:physics:phyx-mini-0656:phyxminiproblems-problemphyxmini0656-massinkilograms, def:physics:phyx-mini-0656:phyxminiproblems-problemphyxmini0656-energyinjoules, def:physics:phyx-mini-0656:phyxminiproblems-problemphyxmini0656-rigidboxelectronsetup}
  Positivity conditions selecting a nondegenerate physical setup
\end{definition}
\begin{proof}
  This is the declared model definition of Has Physical Rigid Box Parameters.
\end{proof}
\begin{definition}[Answer Choice]
  \label{def:physics:phyx-mini-0656:phyxminiproblems-problemphyxmini0656-answerchoice}
  \lean{PhyXMiniProblems.ProblemPhyXMini0656.AnswerChoice}
  Labels of the four displayed answer choices
\end{definition}
\begin{proof}
  This is the declared model definition of Answer Choice.
\end{proof}
\begin{definition}[displayed Length In Nanometers]
  \label{def:physics:phyx-mini-0656:phyxminiproblems-problemphyxmini0656-displayedlengthinnanometers}
  \lean{PhyXMiniProblems.ProblemPhyXMini0656.displayedLengthInNanometers}
  \uses{def:physics:phyx-mini-0656:phyxminiproblems-problemphyxmini0656-answerchoice}
  Nanometre value printed beside each answer label
\end{definition}
\begin{proof}
  This is the declared model definition of displayed Length In Nanometers.
\end{proof}
\begin{definition}[recorded Dataset Answer]
  \label{def:physics:phyx-mini-0656:phyxminiproblems-problemphyxmini0656-recordeddatasetanswer}
  \lean{PhyXMiniProblems.ProblemPhyXMini0656.recordedDatasetAnswer}
  \uses{def:physics:phyx-mini-0656:phyxminiproblems-problemphyxmini0656-answerchoice}
  Answer label recorded in the dataset metadata
\end{definition}
\begin{proof}
  This is the declared model definition of recorded Dataset Answer.
\end{proof}
\begin{definition}[rounded To Nearest Tenth]
  \label{def:physics:phyx-mini-0656:phyxminiproblems-problemphyxmini0656-roundedtonearesttenth}
  \lean{PhyXMiniProblems.ProblemPhyXMini0656.roundedToNearestTenth}
  Round a nanometre readout to the displayed tenth-nanometre precision
\end{definition}
\begin{proof}
  This is the declared model definition of rounded To Nearest Tenth.
\end{proof}
\begin{definition}[Matches Displayed Length]
  \label{def:physics:phyx-mini-0656:phyxminiproblems-problemphyxmini0656-matchesdisplayedlength}
  \lean{PhyXMiniProblems.ProblemPhyXMini0656.MatchesDisplayedLength}
  \uses{def:physics:phyx-mini-0656:phyxminiproblems-problemphyxmini0656-lengthinnanometers, def:physics:phyx-mini-0656:phyxminiproblems-problemphyxmini0656-rigidboxelectronsetup, def:physics:phyx-mini-0656:phyxminiproblems-problemphyxmini0656-answerchoice, def:physics:phyx-mini-0656:phyxminiproblems-problemphyxmini0656-displayedlengthinnanometers, def:physics:phyx-mini-0656:phyxminiproblems-problemphyxmini0656-roundedtonearesttenth}
  A displayed choice agrees with the calculated rounded box length
\end{definition}
\begin{proof}
  This is the declared model definition of Matches Displayed Length.
\end{proof}
\begin{definition}[Is Unique Matching Displayed Length]
  \label{def:physics:phyx-mini-0656:phyxminiproblems-problemphyxmini0656-isuniquematchingdisplayedlength}
  \lean{PhyXMiniProblems.ProblemPhyXMini0656.IsUniqueMatchingDisplayedLength}
  \uses{def:physics:phyx-mini-0656:phyxminiproblems-problemphyxmini0656-rigidboxelectronsetup, def:physics:phyx-mini-0656:phyxminiproblems-problemphyxmini0656-answerchoice, def:physics:phyx-mini-0656:phyxminiproblems-problemphyxmini0656-matchesdisplayedlength}
  The selected label is the unique displayed choice matching the length
\end{definition}
\begin{proof}
  This is the declared model definition of Is Unique Matching Displayed Length.
\end{proof}
\begin{lemma}[mode Number eq four]
  \label{lem:physics:phyx-mini-0656:phyxminiproblems-problemphyxmini0656-modenumber-eq-four}
  \lean{PhyXMiniProblems.ProblemPhyXMini0656.modeNumber_eq_four}
  \uses{def:physics:phyx-mini-0656:phyxminiproblems-problemphyxmini0656-rigidboxelectronsetup, def:physics:phyx-mini-0656:phyxminiproblems-problemphyxmini0656-matchessuppliedwavefunctionfigure, def:physics:phyx-mini-0656:phyxminiproblems-problemphyxmini0656-representsrigidboxstationarymode}
  Four alternating lobes in the supplied graph identify the fourth stationary mode. The numerical mode is derived from raw figure evidence and the general lobe-count interpretation
\end{lemma}
\begin{proof}
  The conclusion follows from the governing relations and figure readouts named in the statement of mode Number eq four.
\end{proof}
\begin{lemma}[box Length within one hundredth nanometer]
  \label{lem:physics:phyx-mini-0656:phyxminiproblems-problemphyxmini0656-boxlength-within-one-hundredth-nanometer}
  \lean{PhyXMiniProblems.ProblemPhyXMini0656.boxLength_within_one_hundredth_nanometer}
  \uses{def:physics:phyx-mini-0656:phyxminiproblems-problemphyxmini0656-lengthinnanometers, def:physics:phyx-mini-0656:phyxminiproblems-problemphyxmini0656-rigidboxelectronsetup, def:physics:phyx-mini-0656:phyxminiproblems-problemphyxmini0656-matchesrigidboxelectronscenario, def:physics:phyx-mini-0656:phyxminiproblems-problemphyxmini0656-matchesproblemenergyreadout, def:physics:phyx-mini-0656:phyxminiproblems-problemphyxmini0656-useselectronmassreferencedata, def:physics:phyx-mini-0656:phyxminiproblems-problemphyxmini0656-matchessuppliedwavefunctionfigure, def:physics:phyx-mini-0656:phyxminiproblems-problemphyxmini0656-representsrigidboxstationarymode, def:physics:phyx-mini-0656:phyxminiproblems-problemphyxmini0656-satisfiesinfiniterigidboxenergylaw, def:physics:phyx-mini-0656:phyxminiproblems-problemphyxmini0656-hasphysicalrigidboxparameters}
  For mode n = 4, energy 6.0 eV, the electron reference mass, and Physlib's ℏ, the spectrum gives approximately 1.001 nm. The stated tolerance is tighter than the tenth-nanometre precision of the answer choices
\end{lemma}
\begin{proof}
  The conclusion follows from the governing relations and figure readouts named in the statement of box Length within one hundredth nanometer.
\end{proof}
% --- Archon declaration topology end ---
% --- Archon physics formalization source end ---
